% GENERATED FILE. Edit canonical lessons, then run npm run generate:books.
% canonical-source-hash: fnv1a64:85c25aaf8372c371
% canonical-lessons: FR-C440-faubourg, FR-C440-arrondissement, FR-C440-prefecture, FR-C440-commissariat, FR-C440-caserne

\chapter{Les quartiers}
\label{ch:fr-les-quartiers}
\hlchaptermodality{440}

\begin{chapteropening}
\textbf{By the end of this chapter:} \emph{I can talk about parts of a town and official buildings: the old suburb, the district, the prefecture, the police station and the barracks.}

Complete the last lesson of chapter 440: I can talk about parts of a town and official buildings: the old suburb, the district, the prefecture, the police station and the barracks.
\end{chapteropening}

\section[le faubourg]{le faubourg — the old suburb, the outer district}

\label{lesson:FR-C440-faubourg}



\begin{quote}
Before the new one: say the French for the belfry, then the French for the town walls.
\end{quote}

\subsection*{You'll want to know: le faubourg}
\textbf{le faubourg} — \textquotedblleft{}the old suburb, the outer district\textquotedblright{}.

\begin{quote}
Ils habitent dans un faubourg tranquille.
\end{quote}

\begin{quote}
\emph{They live in a quiet old suburb.}
\end{quote}

First the part of a town outside its walls.

\begin{cousinweb}
\emph{faubourg} comes from Old French \emph{fors bourg} 'settlement outside the town' (EtymDB), from \emph{fors} 'outside', from Latin \emph{forīs}, and \emph{bourg} 'town'.
\end{cousinweb}

\begin{grammarlens}[title={What you've built}]
One more word for telling, explaining and describing at B1.
\end{grammarlens}

\subsection*{Your turn}
\begin{itemize}
\raggedright
  \item \emph{Say it:} \emph{le faubourg}
  \item \emph{Say it:} \emph{le faubourg}, once more
  \item \emph{Say it:} a sentence of your own with \emph{le faubourg}
  \item \emph{Recall:} say the French for the belfry, then the French for the town walls, then say \emph{le faubourg} again
\end{itemize}

\begin{tcolorbox}[breakable,colback=teal!4,colframe=teal!35!black,title={Before you move on}]
What does le faubourg mean? (\textquotedblleft{}The old suburb, the outer district\textquotedblright{}.) Say it once more.
\end{tcolorbox}
\section[l'arrondissement]{l'arrondissement — the district (of a large city)}

\label{lesson:FR-C440-arrondissement}



\begin{quote}
Before the new one: say the French for the town walls, then the French for the old suburb.
\end{quote}

\subsection*{You'll want to know: l'arrondissement}
\textbf{l'arrondissement} — \textquotedblleft{}the district (of a large city)\textquotedblright{}.

\begin{quote}
Elle travaille dans le cinquième arrondissement.
\end{quote}

\begin{quote}
\emph{She works in the fifth district.}
\end{quote}

A masculine noun. Some large cities are divided into numbered \emph{arrondissements}.

\begin{cousinweb}
\emph{arrondissement} is built on \emph{arrondir} 'to round off', from \emph{rond} 'round', which goes back to Latin \emph{rotundus}, from \emph{rota} 'wheel' (Lewis \& Short).
\end{cousinweb}

\begin{grammarlens}[title={What you've built}]
One more word for telling, explaining and describing at B1.
\end{grammarlens}

\subsection*{Your turn}
\begin{itemize}
\raggedright
  \item \emph{Say it:} \emph{l'arrondissement}
  \item \emph{Say it:} \emph{l'arrondissement}, once more
  \item \emph{Say it:} a sentence of your own with \emph{l'arrondissement}
  \item \emph{Recall:} say the French for the town walls, then the French for the old suburb, then say \emph{l'arrondissement} again
\end{itemize}

\begin{tcolorbox}[breakable,colback=teal!4,colframe=teal!35!black,title={Before you move on}]
What does l'arrondissement mean? (\textquotedblleft{}The district (of a large city)\textquotedblright{}.) Say it once more.
\end{tcolorbox}
\section[la préfecture]{la préfecture — the prefecture (regional state office)}

\label{lesson:FR-C440-prefecture}



\begin{quote}
Before the new one: say the French for the old suburb, then the French for the district.
\end{quote}

\subsection*{You'll want to know: la préfecture}
\textbf{la préfecture} — \textquotedblleft{}the prefecture (regional state office)\textquotedblright{}.

\begin{quote}
Pour son titre de séjour, il va à la préfecture.
\end{quote}

\begin{quote}
\emph{For his residence permit he goes to the prefecture.}
\end{quote}

Where the State's office for a \emph{département} is.

\begin{cousinweb}
\emph{préfecture} is a learned borrowing from Latin \emph{praefectūra} 'office of a prefect', from \emph{praeficere} 'to put in charge', built on \emph{facere} 'to do, to make' (Lewis \& Short).
\end{cousinweb}

\begin{grammarlens}[title={What you've built}]
One more word for telling, explaining and describing at B1.
\end{grammarlens}

\subsection*{Your turn}
\begin{itemize}
\raggedright
  \item \emph{Say it:} \emph{la préfecture}
  \item \emph{Say it:} \emph{la préfecture}, once more
  \item \emph{Say it:} a sentence of your own with \emph{la préfecture}
  \item \emph{Recall:} say the French for the old suburb, then the French for the district, then say \emph{la préfecture} again
\end{itemize}

\begin{tcolorbox}[breakable,colback=teal!4,colframe=teal!35!black,title={Before you move on}]
What does la préfecture mean? (\textquotedblleft{}The prefecture (regional state office)\textquotedblright{}.) Say it once more.
\end{tcolorbox}
\section[le commissariat]{le commissariat — the police station}

\label{lesson:FR-C440-commissariat}



\begin{quote}
Before the new one: say the French for the district, then the French for the prefecture.
\end{quote}

\subsection*{You'll want to know: le commissariat}
\textbf{le commissariat} — \textquotedblleft{}the police station\textquotedblright{}.

\begin{quote}
Après le vol de son sac, elle est allée au commissariat.
\end{quote}

\begin{quote}
\emph{After her bag was stolen, she went to the police station.}
\end{quote}

In full: \emph{le commissariat de police}.

\begin{cousinweb}
\emph{commissariat} is built on \emph{commissaire} 'commissioner', from Medieval Latin \emph{commissārius}, from \emph{committere} 'to entrust', built on \emph{mittere} 'to send' (Lewis \& Short).
\end{cousinweb}

\begin{grammarlens}[title={What you've built}]
One more word for telling, explaining and describing at B1.
\end{grammarlens}

\subsection*{Your turn}
\begin{itemize}
\raggedright
  \item \emph{Say it:} \emph{le commissariat}
  \item \emph{Say it:} \emph{le commissariat}, once more
  \item \emph{Say it:} a sentence of your own with \emph{le commissariat}
  \item \emph{Recall:} say the French for the district, then the French for the prefecture, then say \emph{le commissariat} again
\end{itemize}

\begin{tcolorbox}[breakable,colback=teal!4,colframe=teal!35!black,title={Before you move on}]
What does le commissariat mean? (\textquotedblleft{}The police station\textquotedblright{}.) Say it once more.
\end{tcolorbox}
\section[la caserne]{la caserne — the barracks; the fire station}

\label{lesson:FR-C440-caserne}



\begin{quote}
Before the new one: say the French for the prefecture, then the French for the police station.
\end{quote}

\subsection*{You'll want to know: la caserne}
\textbf{la caserne} — \textquotedblleft{}the barracks; the fire station\textquotedblright{}.

\begin{quote}
La caserne des pompiers est à côté de l'école.
\end{quote}

\begin{quote}
\emph{The fire station is next to the school.}
\end{quote}

\emph{La caserne des pompiers} — the fire station.

\begin{cousinweb}
\emph{caserne} was borrowed from Provençal \emph{cazerna} 'small hut for four soldiers', which goes back to Latin \emph{quaterna} 'group of four' (EtymDB).
\end{cousinweb}

\begin{grammarlens}[title={What you've built}]
One more word for telling, explaining and describing at B1.
\end{grammarlens}

\subsection*{Your turn}
\begin{itemize}
\raggedright
  \item \emph{Say it:} \emph{la caserne}
  \item \emph{Say it:} \emph{la caserne}, once more
  \item \emph{Say it:} a sentence of your own with \emph{la caserne}
  \item \emph{Recall:} say the French for the prefecture, then the French for the police station, then say \emph{la caserne} again
\end{itemize}

\begin{tcolorbox}[breakable,colback=teal!4,colframe=teal!35!black,title={Before you move on}]
What does la caserne mean? (\textquotedblleft{}The barracks; the fire station\textquotedblright{}.) Say it once more.
\end{tcolorbox}
